% !TeX root = ../../portafoliofinal.tex
% !TeX encoding = utf8
% ====================================================================
%  BIBLIOGRAFÍA UNIFICADA — PORTAFOLIO FINAL
%  Fuentes consultadas para las tres prácticas y el informe global
% ====================================================================

\chapter*{Fuentes consultadas}
\addcontentsline{toc}{chapter}{Fuentes consultadas}

A continuación se recogen, de forma unificada, todas las fuentes consultadas
para la elaboración del portafolio grupal final, organizadas por ámbito temático.

% ================================================================
\section*{Fuentes corporativas y documentos oficiales de CaixaBank}
% ================================================================

\begin{enumerate}[label={[\arabic*]}]

  \item \textbf{CaixaBank} --- \textit{Plan Estratégico 2025--2027 «Uno a Uno»}.
        Presentación ante inversores (\textit{Investor Day}), noviembre de 2024.\\
        \url{https://www.caixabank.com/es/sobre-caixabank/nuestro-plan-estrategico.html}

  \item \textbf{CaixaBank} --- \textit{Informe Anual 2024}. Datos corporativos:
        plantilla consolidada, red de oficinas y cifras clave de gestión del capital humano.\\
        \url{https://www.caixabank.com/informacionparaaccionistaseinversores/informacioneconomicofinanciera/informeanual_es.html}

  \item \textbf{CaixaBank} --- \textit{Informe de Sostenibilidad e Información
        No Financiera 2023}. Datos sobre estructura salarial, brecha de género,
        benchmarking retributivo y programas de bienestar.\\
        \url{https://www.caixabank.com/es/investor-relations/annual-reports.html}

  \item \textbf{CaixaBank} --- \textit{Política de Retribución}.
        Documento aprobado por el Consejo de Administración con criterios de equidad,
        retribución variable e incentivos a largo plazo.\\
        \url{https://www.caixabank.com/es/investor-relations/corporate-governance/remuneration-policy.html}

  \item \textbf{CaixaBank} --- \textit{Nota de prensa: Integración con Bankia} (2021).
        Información sobre el proceso de fusión y la gestión del capital humano derivada
        de la unión de ambas entidades.\\
        \url{https://www.caixabank.com/es/prensa/2021/caixabank-bankia-integracion.html}

  \item \textbf{CaixaBank} --- \textit{Adelanto del pago de prestaciones ERTE} (abril de 2020).
        CaixaBank adelantó el pago de 100 millones de euros en prestaciones por desempleo
        para 138.000 personas afectadas por ERTE.\\
        \url{https://www.caixabank.com/es/actualidad/noticias/caixabank-adelanta-el-pago-de-100-millones-de-euros-en-prestaciones-por-desempleo-para-138-000-personas-afectadas-por-erte}

  \item \textbf{CaixaBank --- Sostenibilidad: Personas}.
        Política de formación, programas de talento, diversidad y bienestar de la plantilla.\\
        \url{https://www.caixabank.com/es/sostenibilidad/personas.html}

\end{enumerate}

% ================================================================
\section*{Portales de empleo y selección de CaixaBank}
% ================================================================

\begin{enumerate}[label={[\arabic*]}, resume]

  \item \textbf{CaixaBankCareers} --- Portal oficial de empleo.\\
        \url{https://caixabankcareers.com/}

  \item \textbf{CaixaBankCareers} --- Proceso de selección.\\
        \url{https://caixabankcareers.com/content/Proceso-de-seleccion/?locale=es_ES}

  \item \textbf{CaixaBankCareers} --- ¿Qué te ofrecemos? Desarrollo profesional,
        programas de talento y plan de movilidad interna (\textit{moving}).\\
        \url{https://caixabankcareers.com/content/Que-te-ofrecemos/?locale=es_ES}

  \item \textbf{LinkedIn --- CaixaBank}. Perfil corporativo; estrategia de
        \textit{employer branding} y publicación de ofertas (reclutamiento 3.0).\\
        \url{https://es.linkedin.com/company/caixabank/jobs}

\end{enumerate}

% ================================================================
\section*{Fuentes sobre el ERE CaixaBank 2021}
% ================================================================

\begin{enumerate}[label={[\arabic*]}, resume]

  \item \textbf{Cinco Días (El País)} --- «CaixaBank y sindicatos acuerdan el ERE
        con 6.452 salidas y un coste de 1.900 millones de euros» (1 jul. 2021).\\
        \url{https://cincodias.elpais.com/cincodias/2021/07/01/companias/1625130605_136870.html}

  \item \textbf{El Economista} --- «CaixaBank sella el ERE con los sindicatos
        con 6.452 salidas, 1.839 menos, todas ellas voluntarias» (jul. 2021).\\
        \url{https://www.eleconomista.es/empresas-finanzas/noticias/11303640/07/21}

  \item \textbf{elDiario.es} --- «CaixaBank y los sindicatos cierran el mayor
        ERE bancario con 6.452 salidas» (1 jul. 2021).\\
        \url{https://www.eldiario.es/economia/caixabank-sindicatos-cierran-mayor-ere-bancario-6-452-salidas_1_8102010.html}

\end{enumerate}

% ================================================================
\section*{Tecnología y SIRH}
% ================================================================

\begin{enumerate}[label={[\arabic*]}, resume]

  \item \textbf{SAP} --- CaixaBank \& SAP SuccessFactors. Implantación de la
        plataforma de RRHH en la nube: gestión del aprendizaje, evaluación del
        desempeño, planificación de la sucesión y People Analytics.\\
        \url{https://www.sap.com/spain/customer-testimonials/caixabank.html}

  \item \textbf{Microsoft} --- «CaixaBank adopta Microsoft 365 Copilot».
        Integración de IA generativa para productividad y personalización del
        aprendizaje en la plataforma Evoluciona (antes denominada Virtaula).\\
        \url{https://news.microsoft.com/es-es/2024/caixabank-microsoft-copilot/}

  \item \textbf{People Xperience Hub / Dialnet} --- Artículo académico sobre
        la iniciativa RPO de CaixaBank (2019): estandarización y digitalización
        del ciclo de atracción de talento.\\
        \url{https://dialnet.unirioja.es/servlet/articulo?codigo=7731413}

\end{enumerate}

% ================================================================
\section*{Fuentes sobre desarrollo profesional y formación}
% ================================================================

\begin{enumerate}[label={[\arabic*]}, resume]

  \item \textbf{Whitmore, John} --- \textit{Coaching for Performance: GROWing
        Human Potential and Purpose} (5ª ed.). Nicholas Brealey Publishing,
        Londres, 2017. Referencia para el método GROW y el \textit{shadow coaching}.

  \item \textbf{Harvard Business Review} --- «Why Reverse Mentoring Works and
        How to Do It Right». HBR, octubre de 2019.\\
        \url{https://hbr.org/2019/10/why-reverse-mentoring-works-and-how-to-do-it-right}

  \item \textbf{Gallup} --- «3 Reasons Why Performance Development Wins in the
        Workplace». Fundamento para la evaluación continua y el feedback frecuente.\\
        \url{https://www.gallup.com/workplace/231620/why-performance-development-wins-workplace.aspx}

  \item \textbf{Gallup} --- «Is Your Employee Recognition Really Authentic?».
        Sobre reconocimiento auténtico y mejora del compromiso.\\
        \url{https://www.gallup.com/workplace/508208/employee-recognition-really-authentic.aspx}

  \item \textbf{AEDIPE} --- Artículos sobre \textit{outdoor training},
        gamificación en RRHH y tendencias en formación corporativa en el
        sector financiero español.\\
        \url{https://www.aedipe.es/}

  \item \textbf{Llorente \& Cuenca / Equipos \& Talento} --- «La gamificación
        en la formación corporativa: casos de éxito en el sector bancario
        español».\\
        \url{https://www.equiposytalento.com/}

  \item \textbf{Muñiz, R. / Marketing XXI} --- «El outdoor training como
        herramienta de desarrollo directivo».\\
        \url{https://www.marketing-xxi.com/outdoor-training.html}

\end{enumerate}

% ================================================================
\section*{Fuentes sobre retribución y convenio colectivo}
% ================================================================

\begin{enumerate}[label={[\arabic*]}, resume]

  \item \textbf{IV Convenio Colectivo del Sector de Banca} --- Negociado entre
        la AEB, CCOO, UGT y FINE. Categorías profesionales y tablas salariales
        mínimas. Publicado en el BOE.\\
        \url{https://www.boe.es/diario_boe/txt.php?id=BOE-A-2023-14936} % chktex 8

  \item \textbf{Glassdoor / Indeed} --- Reseñas de empleados de CaixaBank:
        retribución, beneficios y medidas de conciliación.\\
        \url{https://www.glassdoor.es/Resenas/CaixaBank-resenas-E17906.htm}

  \item \textbf{Expansión / Cinco Días} --- Artículos sobre la política salarial
        del sector bancario español y el posicionamiento retributivo de CaixaBank.\\
        \url{https://www.expansion.com/}

\end{enumerate}

% ================================================================
\section*{Fuentes primarias (entrevistas)}
% ================================================================

\begin{enumerate}[label={[\arabic*]}, resume]

  \item \textbf{Entrevista 1} --- Director de Recursos Humanos regional de CaixaBank
        (anónimo a petición). Entrevista telefónica, 13 de marzo de 2026,
        duración aproximada 15 minutos. Temática: Análisis de puestos y
        Planificación de RRHH (Temas 2 y 3). Véase Apéndice~A.

  \item \textbf{Entrevista 2} --- Director de Recursos Humanos regional de CaixaBank
        (anónimo a petición). Entrevista telefónica, 12 de abril de 2026,
        duración aproximada 15 minutos. Temática: Reclutamiento, Selección
        e Integración (Temas 4 y 5). Véase Apéndice~B.

  \item \textbf{Entrevista 3} --- Directora de sucursal de CaixaBank
        (anónima a petición). Entrevista semiestructurada, 11 de abril de 2026.
        Temática: Reclutamiento, selección, herramientas digitales y videoentrevistas.

\end{enumerate}

% ================================================================
\section*{Documentos teóricos de la asignatura}
% ================================================================

\begin{enumerate}[label={[\arabic*]}, resume]

  \item Apuntes de clase de los Temas 2, 3, 4, 5, 6, 7 y 8.

  \item \textbf{Guiones para el desarrollo de las prácticas grupales},
        Dirección de Recursos Humanos~I, Doble Grado ADE-Ingenierías,
        Universidad de Granada, curso 2025--2026.

  \item \textbf{Instrucciones detalladas para la elaboración del portafolio
        grupal}, Javier Aguilera Caracuel, Universidad de Granada, mayo de 2025.

\end{enumerate}

\endinput
